\section*{第 \ref{ch:b2:blood-pressure} 章——血压调节与运动}

\begin{solution}{exo:b2:blood-pressure:1}
$\bar P - P_{\text{静脉}} = Q\,R = f\,V_s\,R$。心率（
\omterm{prop:b2:heart:conduction}{窦房结}，受迷走神经和交感神经支配）；\omterm{prop:b2:heart:cycle}{每搏输出量}（
\omterm{def:b2:heart:anatomy}{心室}：其充盈由\omterm{def:b2:blood-circulation:vessels}{静脉}决定，其收缩力由交感神经决定）；阻力（
\omterm{def:b2:blood-circulation:vessels}{小动脉}平滑肌，受交感神经、激素
和局部代谢产物调节
）。
\end{solution}

\begin{solution}{exo:b2:blood-pressure:2}
感受器：颈\omterm{def:b2:blood-circulation:vessels}{动脉}窦和主\omterm{def:b2:blood-circulation:vessels}{动脉}弓中的牵张感受器。
中枢：延髓的心血管中枢。效应器：
迷走神经（心率）和交感神经（心率、收缩力，
以及\omterm{def:b2:blood-circulation:vessels}{小动脉}和\omterm{def:b2:blood-circulation:vessels}{静脉}张力）。负反馈，延迟约一秒。
若无感受器，一天的平均血压仍保持正常，
但每分钟的血压会剧烈摆动，
随体位和情绪而变化。
\end{solution}

\begin{solution}{exo:b2:blood-pressure:3}
肾脏（\omterm{def:b2:blood-circulation:vessels}{小动脉}细胞）在灌注或钠输送下降、
或交感神经兴奋时释放肾素；肾素裂解
血管紧张素原（肝脏）生成血管紧张素 I；肺中的酶
将其转化为血管紧张素 II，后者使\omterm{def:b2:blood-circulation:vessels}{小动脉}收缩，
并刺激肾上腺皮质分泌\omterm{prop:b2:blood-pressure:raas}{醛固酮}；\omterm{prop:b2:blood-pressure:raas}{醛固酮}使肾脏
潴留钠和水，从而提高血容量。
\end{solution}

\begin{solution}{exo:b2:blood-pressure:4}
肌肉自身的\omterm{def:b2:blood-circulation:vessels}{小动脉}因代谢产物而舒张；
\omterm{thm:b2:heart:starling}{心输出量}增至五倍，由交感神经和肌泵驱动；
肠道、肾脏和静息肌肉的\omterm{def:b2:blood-circulation:vessels}{小动脉}受交感神经
收缩，因而让出
血流。
\end{solution}

\begin{solution}{exo:b2:blood-pressure:5}
$83\times 1.2 = \qty{100}{mmHg}$; $333\times 0.35 = \qty{117}{mmHg}$.
阻力下降 $1.2/0.35 = 3.4$ 倍：半径增大为 $3.4^{1/4} =
1.36$ 倍。
\end{solution}

\begin{solution}{exo:b2:blood-pressure:6}
$30/(1 + G)$：分别为 10、6 和 \qty{3.3}{mmHg}。增益减半后，
残余下降量加倍：她站起时头晕、眼前发灰，
可能昏倒。
\end{solution}

\begin{solution}{exo:b2:blood-pressure:7}
$\rho g h$: $1060\times 9.81\times 0.4 = \qty{4160}{Pa} =
\qty{31}{mmHg}$; $1060\times 9.81\times 1.2 = \qty{12500}{Pa} =
\qty{94}{mmHg}$。脑处 $95 - 31 = \qty{64}{mmHg}$；踝处 $95 + 94 =
\qty{189}{mmHg}$。长颈鹿头部位于 \qty{2.5}{m} 高处，需克服 \qty{195}{mmHg}：
其\omterm{def:b2:heart:anatomy}{心脏}处的平均血压须达约 \qty{250}{mmHg}，并且
腿部血管壁厚、皮肤紧实。
\end{solution}

\begin{solution}{exo:b2:blood-pressure:8}
$Q = 280/(190 - 130) = \qty{4.7}{L/min}$。运动员：$5000/(200 - 30) =
\qty{29}{L/min}$；\omterm{prop:b2:heart:cycle}{每搏输出量} $\num{29400}/190 = \qty{155}{mL}$。
\end{solution}

\begin{solution}{exo:b2:blood-pressure:9}
$35/5 = \qty{7}{mmHg}$ 的下降，血压约为 \qty{88}{mmHg}。恢复
血容量的途径有：组织液重吸收进入\omterm{def:b2:blood-circulation:vessels}{毛细血管}
（数分钟至一小时）；在\omterm{prop:b2:blood-pressure:raas}{抗利尿激素}和\omterm{prop:b2:blood-pressure:raas}{醛固酮}作用下尿量减少
（数小时至数天）；口渴与饮水（数小时）；\omterm{def:b2:blood-circulation:blood}{红细胞}来自骨髓（数周）。
皮肤苍白、发凉，是因为反射使其\omterm{def:b2:blood-circulation:vessels}{小动脉}收缩，
以便把血压留给大脑
和\omterm{def:b2:heart:anatomy}{心脏}。
\end{solution}

\begin{solution}{exo:b2:blood-pressure:10}
肌肉在 $R = 0.5$ 时：$95/0.5 = \qty{190}{mL/s}$，即 \qty{11.4}{L/min}；
若加上其他器官的 \qty{4}{L/min}，\omterm{def:b2:heart:anatomy}{心脏}需输出
\qty{15.4}{L/min}。若 $Q$ 固定为 \qty{5}{L/min}，总阻力由 1.14 降至
$1/(1/1.14 - 1/5 + 1/0.5) =
\qty{0.37}{mmHg.s/mL}$，血压降至 \qty{31}{mmHg}——大脑将昏厥。
机体会将输出量向第一种情形靠拢，
并使其他血管床收缩，从而维持血压，
同时保证肌肉获得血流。
\end{solution}

\begin{solution}{exo:b2:blood-pressure:11}
(a) 灌注不足的肾脏感知到低血压，释放肾素，
血管紧张素和\omterm{prop:b2:blood-pressure:raas}{醛固酮}使阻力和血容量增加。(b) \omterm{def:b2:blood-pressure:baroreflex}{压力反射}能缓冲变化，
但会重设到持续存在的水平，
无法改变肾脏对血容量的潴留。
(c) 肾素的来源，以及需要更高血压的器官，
均已不复存在。(d) 阻断了向血管紧张素 II 的转化：
阻力降低，\omterm{prop:b2:blood-pressure:raas}{醛固酮}减少，血压下降。
\end{solution}

\begin{solution}{exo:b2:blood-pressure:12}
\omterm{def:b2:blood-pressure:baroreflex}{压力反射}在一秒内起作用，增益约为四，调定点
随血压漂移：它能抑制波动，但无法维持某一水平。
肾脏则在数小时至数天内起作用，实际上增益无限大——
它持续排泄或潴留，直至达到所设定的血压——
其调定点由自身生理决定，
因而能维持某一水平。因此，慢性血压升高意味着肾脏的调定点已经移动，
只有作用于肾脏调节环路，
或作用于其所需阻力的药物，才能使血压降下来。
\end{solution}

\begin{solution}{pb:b2:blood-pressure:1}
\textbf{1.} $1/R = 1/7.6 + 1/22.8 + 1/4.1 + 1/5.2 + 1/5.7 + 1/19 +
1/28.5 = 0.875$: $R = \qty{1.14}{mmHg.s/mL}$; $\bar P/Q = 95/83.3 =
1.14$.
\textbf{2.} $\bar P/R_i$：脑 0.75，\omterm{def:b2:heart:anatomy}{心脏} 0.25，肠道和肝脏 1.39，
肾脏 1.10，肌肉 1.0，皮肤 0.30，其他 \qty{0.20}{L/min}；
所占比例分别为 15、5、28、22、20、6 和 \qty{4}{\%}。
\textbf{3.} $5000/70 = \qty{71}{mL}$；每分钟
$5\times 50 = \qty{250}{mL}$ 的氧气。
\textbf{4.} 脑 $0.75/1.4 = \qty{0.54}{L/min/kg}$；肾脏 $1.1/0.3
= \qty{3.7}{L/min/kg}$，是脑的七倍，是全身平均值的五十倍：
肾脏被灌注是为了过滤，而非为了供能。
\textbf{5.} 所有器官共用同一血压，各器官所获血流
由其自身阻力决定；若集中调节血流，
则任何需求发生变化的器官都会缺血。血压是
共同的“货币”。
\textbf{6.} \omterm{def:b2:heart:anatomy}{心脏}：$250\times 0.2\times 0.7 = \qty{35}{mL/min}$；
肌肉：$1000\times 0.2\times 0.25 = \qty{50}{mL/min}$。\omterm{def:b2:heart:anatomy}{心脏}
已摄取了流经\omterm{def:b2:blood-circulation:blood}{血液}中的大部分氧，因此只能靠增加血流
获得更多氧，这正是其工作加强时\omterm{def:b2:blood-circulation:vessels}{小动脉}立即舒张的
原因。
\textbf{7.} $Q$ 下降 \qty{30}{\%}（$R$ 不变）：$\bar P$ 下降
\qty{30}{\%}，即 \qty{28.5}{mmHg}，降至 \qty{66.5}{mmHg}。
\textbf{8.} $28.5/5 = \qty{5.7}{mmHg}$：约为 \qty{89}{mmHg}。
\textbf{9.} $Q = 85\times 0.7\times 71 = \qty{4.2}{L/min} =
\qty{70}{mL/s}$；$R = 89.3/70 = \qty{1.27}{mmHg.s/mL}$，增加
\qty{11}{\%}.
\textbf{10.} $1060\times 9.81\times 0.4 = \qty{4160}{Pa} =
\qty{31}{mmHg}$；反射前脑部 $66.5 - 31 = \qty{35}{mmHg}$，
反射后为 $89 - 31 = \qty{58}{mmHg}$。
\textbf{11.} $28.5/2 = \qty{14}{mmHg}$：\omterm{def:b2:heart:anatomy}{心脏}处为 \qty{81}{mmHg}，
脑部为 \qty{50}{mmHg}——头晕、眼前发灰，
每次从椅子上站起时都几近昏厥。
\textbf{12.} 平卧消除了 \qty{31}{mmHg} 的静水压损失，
并使淤积的\omterm{def:b2:blood-circulation:blood}{血液}回流至\omterm{def:b2:heart:anatomy}{心脏}；\omterm{prop:b2:heart:cycle}{每搏输出量}和
脑部血压在数次心跳内即可恢复。
\textbf{13.} $45/5 = \qty{9}{mmHg}$: \qty{86}{mmHg}.
\textbf{14.} $1/R = 0.875 - 1/19 - 1/4.1 + 1/38 + 1/8.2 = 0.727$: $R =
\qty{1.38}{mmHg.s/mL}$。脑 $86/7.6 = \qty{11.3}{mL/s} =
\qty{0.68}{L/min}$；皮肤 $86/38 = \qty{2.3}{mL/s} = \qty{0.14}{L/min}$；
肠道 $86/8.2 = \qty{10.5}{mL/s} = \qty{0.63}{L/min}$。
\textbf{15.} 血量减少、\omterm{def:b2:blood-circulation:vessels}{小动脉}收缩后，\omterm{def:b2:blood-circulation:vessels}{毛细血管}
压力沿整条\omterm{def:b2:blood-circulation:vessels}{毛细血管}都低于\omterm{thm:b2:blood-circulation:oncotic}{胶体渗透压}，
因此液体被重吸收而非滤出；\omterm{def:b2:blood-circulation:blood}{血浆}被稀释，
\omterm{def:b2:blood-circulation:blood}{血细胞比容}下降。
\textbf{16.} \omterm{prop:b2:blood-pressure:raas}{醛固酮}经肾素和血管紧张素途径，受
肾脏灌注不足和交感神经驱动所触发；\omterm{prop:b2:blood-pressure:raas}{抗利尿激素}
则由血容量下降和\omterm{def:b2:blood-circulation:blood}{血浆}
浓度升高所触发。
\textbf{17.} $5\times 10^{12}$ 个细胞，按每秒 $2.5\times 10^{6}$ 个计：
$2\times 10^{6}$ 秒，约三周。
\textbf{18.} 血压低于约 \qty{60}{mmHg} 时，反射曲线趋于平坦：
血管已尽可能收缩，心率已尽可能加快，增益为零，
每一份进一步的损失都使血压等量下降；
灌注不足的组织随之释放酸性物质并舒张，
血压随之骤降。
\textbf{19.} 肌肉 $110/0.33 = \qty{333}{mL/s} = \qty{20}{L/min}$；
皮肤 $110/5 = \qty{22}{mL/s} = \qty{1.3}{L/min}$。
\textbf{20.} 脑阻力 $8.8$（110 时血流为 0.75），\omterm{def:b2:heart:anatomy}{心脏} $6.6$
（血流 1.0）：$1/R = 1/8.8 + 1/6.6 + 1/8.2 + 1/10.4 + 1/0.33 + 1/5 +
1/28.5 = 3.75$: $R = \qty{0.27}{mmHg.s/mL}$; $Q = 110/0.27 =
\qty{410}{mL/s} = \qty{25}{L/min}$.
\textbf{21.} $\num{24700}/180 = \qty{137}{mL}$.
\textbf{22.} $24.7\times(200 - 40) = \qty{3950}{mL/min}$，
为静息时 \qty{250}{mL/min} 的十六倍。
\textbf{23.} 血流 $\times 4.9$，提取率 $160/50 = \times 3.2$；
$4.9\times 3.2 = 16$.
\textbf{24.} 来自运动皮层的指令，以及肌肉感受器发出的信号，
重新调整延髓中\omterm{def:b2:blood-pressure:baroreflex}{压力感受器}的输入权重，使反射维持
\qty{110}{mmHg} 而非 95——
即提高调定点，而非关闭环路。
\textbf{25.} 站立时残余下降 \qty{5.7}{mmHg}；失血后为
\qty{86}{mmHg}；运动时为 \qty{25}{L/min}，每分钟摄取
氧气 \qty{4}{L}。
\end{solution}
